\begin{trapbox}
\begin{description}[style=nextline, leftmargin=2.8em, labelsep=0.5em, itemsep=0.35em]
    \item[\textbf{\textcolor{CTrap}{易错点一（切线方程误写）}}] 写错椭圆在点 $P(x_0,y_0)$ 处的切线方程，误用 $xx_0/a^2+yy_0/b^2=1$ 或忘记除以 $a^2,b^2$。
    \item[\textbf{\textcolor{CTrap}{正确理解（标准切线公式）：}}] 正确切线方程为 $\frac{x_0 x}{a^2}+\frac{y_0 y}{b^2}=1$，其中 $(x_0,y_0)$ 在椭圆上。
    \item[\textbf{\textcolor{CTrap}{错因分析（公式记忆混淆）：}}] 将椭圆切线公式与圆切线公式 $xx_0+yy_0=r^2$ 混淆，或未区分 $a^2,b^2$ 的位置。
\end{description}

\medskip
\begin{description}[style=nextline, leftmargin=2.8em, labelsep=0.5em, itemsep=0.35em]
    \item[\textbf{\textcolor{CTrap}{易错点二（距离符号忽略）}}] 计算 $|P_1A_1|$ 时，未考虑 $y$ 坐标的符号，直接代入表达式而未取绝对值。
    \item[\textbf{\textcolor{CTrap}{正确理解（非负距离）：}}] $|P_1A_1|$ 是线段长度，应为非负值，由交点纵坐标与 $A_1$ 纵坐标差的绝对值给出。
    \item[\textbf{\textcolor{CTrap}{错因分析（几何意义不清）：}}] 忽略距离定义，误将坐标差直接作为距离，导致乘积计算出现负号或错误。
\end{description}

\medskip
\begin{description}[style=nextline, leftmargin=2.8em, labelsep=0.5em, itemsep=0.35em]
    \item[\textbf{\textcolor{CTrap}{易错点三（最值条件遗漏）}}] 求面积最小值时，未考虑 $y_0$ 的取值范围，或误以为最小值在 $y_0=0$ 时取得。
    \item[\textbf{\textcolor{CTrap}{正确理解（取等条件）：}}] 面积最小值在 $|y_0|=b$ 即 $P$ 在短轴端点时取得，此时 $x_0=0$。
    \item[\textbf{\textcolor{CTrap}{错因分析（范围分析不当）：}}] 未利用椭圆方程中 $y_0^2 \le b^2$ 的条件，或错误认为 $y_0$ 可趋于 $0$ 导致面积趋于无穷。
\end{description}
\end{trapbox}
